\begin{thebibliography}{99}
\bibitem{GrothCoates1966}
A. Grothendieck, letter to J. Coates dated 6.1.1966 as written,
cote 16, scans 048--050; especially scans 049--050.
CSG transcription L11d, printed p.~3, Proposition.
\url{https://csg.igrothendieck.org/wp-content/uploads/2024/04/L11d.pdf}.

\bibitem{GrothCoates1967}
A. Grothendieck, letter to J. Coates, 4.1.1967,
cote 16 scans 044--047. CSG transcription L10d,
printed p.~2, Proposition (i), (ii bis), and flag corollary.
\url{https://csg.igrothendieck.org/wp-content/uploads/2024/04/L10d.pdf}.

\bibitem{Archive}
A. Grothendieck, Montpellier archive, cotes 16 and 44,
supplied transcriptions numbered by scan.
Cote 16: scans 014--018, 034--040, 049--050.
Cote 44: scan 004, Theorem 5, manuscript p.~2;
scans 008--009 for incidence duality and primitive signs.
\url{https://grothendieck.umontpellier.fr/16.pdf}.

\bibitem{Kleiman1968}
S. L. Kleiman, \emph{Algebraic cycles and the Weil conjectures},
in \emph{Dix expos\'es sur la cohomologie des sch\'emas},
North-Holland/Masson, 1968, pp.~359--386.
Exact locators: 1.4.7; 2.3--2.5; 2.7; 2.12;
Theorem 2.13(i),(ii$'$),(iv), pp.~372--373;
Theorem 2A11(2), p.~377;
Proposition 3.6 and Corollary 3.9, pp.~379--380;
Theorem 3.11 and Corollaries 3.12, 3.14, pp.~380--382;
Lemma 4.6 and Theorem 4.7, pp.~385--386.
\url{https://download.uni-mainz.de/mathematik/Algebraische%20Geometrie/Lehre/WS23.Padische.1966.kleiman.weil.pdf}.

\bibitem{Milne1999a}
J. S. Milne, \emph{Lefschetz classes on abelian varieties},
Duke Mathematical Journal \textbf{96} (1999), 639--675,
Theorem 5.9 (author PDF p.~26), read with the author's correction
to its proof.
\url{https://www.jmilne.org/math/articles/1999aP.pdf};
\url{https://www.jmilne.org/math/articles/1999a.html}.

\bibitem{Milne2002}
J. S. Milne, \emph{Polarizations and Grothendieck's standard
conjectures}, Annals of Mathematics (2) \textbf{155} (2002), 599--610.
Later author version (September 19, 2001):
Theorem 3.3(a),(b), pp.~8--9; Corollary 3.4;
Remarks 3.7--3.8, p.~9.
\url{https://www.jmilne.org/math/articles/2002aP.pdf}.
Original author version (March 26/August 14, 2001):
Theorem 4.6(a),(b), pp.~13--14.
\url{https://www.jmilne.org/math/articles/2002aS.pdf}.
Addendum to Theorem 2.1:
\url{https://www.jmilne.org/math/articles/2002a.html}.

\bibitem{CMHodge032}
OpenAI, \emph{The rational Hodge conjecture for CM abelian
varieties}, October 6, 2026, supplied file 032.
Theorem 1.1, p.~1 and concluding proof p.~47;
Corollary 8.4, p.~50. Used as stipulated cited input.

\bibitem{VanHoftenSempliner}
P. van Hoften and J. Sempliner,
\emph{The Hodge conjecture for CM abelian varieties},
supplied Hexagon 2610.00259v2, Theorem I and \S1.1.2, p.~1.

\bibitem{MilneCM}
J. S. Milne, \emph{Complex Multiplication}, author course notes,
Propositions 3.12--3.13, p.~30.
\url{https://www.jmilne.org/math/CourseNotes/CM.pdf}.

\bibitem{MilneAV}
J. S. Milne, \emph{Abelian Varieties}, course notes,
version 1.1, July 27, 1998, Proposition 13.1, printed p.~53;
for exterior cohomology, Theorem 11.1(b).
The exact version consulted is
\url{https://www.jmilne.org/math/CourseNotes/AV110.pdf}.

\bibitem{Jannsen1992}
U. Jannsen, \emph{Motives, numerical equivalence, and
semi-simplicity}, Inventiones Mathematicae \textbf{107}
(1992), 447--452, Theorem 1.
\url{https://epub.uni-regensburg.de/26642/1/jannsen10.pdf}.

\bibitem{Alpoge}
L. Alp\"oge, \emph{Remark}, supplied four-page note and Lean archive,
20 November 2026. Steps 2--5, 7--10, 13, 17 as numbered in the
supplied note; the definitions and theorem declarations actually
used are listed in the integral sections and manifests.
\url{https://alpo.ge/remark.pdf}.
The supplied \path{ALPOGE_LEAN_MAP_20261120.md} and
\path{ALPOGE_LEAN_SCAN.tsv} are navigation/audit inputs,
not new certificates produced by this run.

\bibitem{HodgeSign}
\emph{HODGE\_SIGN}, supplied
\path{RETURN_HODGE_SIGN_final_20261121T0057Z/HODGE_SIGN.tex}.
Used locators: SG-02; SG-04 with its scope remark;
SG-08--SG-11; SL-04, SL-06--SL-09;
AR-01; U-01--U-05 and Q6g.
GCT\_NSRH F03 and the supplied independent check are cited
as directed by the current prompt.

\bibitem{WeilI}
P. Deligne, \emph{La conjecture de Weil. I},
Publications Math\'ematiques de l'IH\'ES \textbf{43}
(1974), 273--307.
Theorem 1.6; (4.1), (4.3.3); (5.4), (5.8)(a),(b).
\url{https://www.numdam.org/item/PMIHES_1974__43__273_0.pdf}.

\bibitem{WeilII}
P. Deligne, \emph{La conjecture de Weil. II},
Publications Math\'ematiques de l'IH\'ES \textbf{52}
(1980), 137--252.
Theorem 3.3.1 over a point; Theorem 4.1.1.
\url{https://www.numdam.org/item/PMIHES_1980__52__137_0.pdf}.

\bibitem{SGA4}
M. Artin, A. Grothendieck and J.-L. Verdier (eds.),
\emph{Th\'eorie des topos et cohomologie \'etale des sch\'emas},
SGA 4, Lecture Notes in Mathematics 269, 270, 305.
Expos\'e XVI, Corollary 2.2, with locally constant coefficients;
Expos\'e XVIII for duality.

\bibitem{SGA5}
A. Grothendieck, edited by L. Illusie,
\emph{Cohomologie $\ell$-adique et fonctions $L$},
SGA 5, Lecture Notes in Mathematics 589, 1977.
Expos\'e III, \S3, 5.1--5.3;
Expos\'e VII, Theorem 7.1 and Corollary 7.2.

\bibitem{SGA6}
P. Berthelot, A. Grothendieck and L. Illusie,
\emph{Th\'eorie des intersections et th\'eor\`eme de
Riemann--Roch}, SGA 6, Lecture Notes in Mathematics 225, 1971.
Expos\'e V, Definition 2.1, Example 2.2(b), Lemma 2.2.1.

\bibitem{SGA7}
A. Grothendieck, P. Deligne and N. Katz (eds.),
\emph{Groupes de monodromie en g\'eom\'etrie alg\'ebrique},
SGA 7 I--II, Lecture Notes in Mathematics 288 and 340.
I, Expos\'e IX, Corollary 2.2.9 and Proposition 3.2,
3.2.1; II, Expos\'e XIV, 3.2,
and Expos\'e XVIII, 6.6.

\bibitem{EGA}
A. Grothendieck and J. Dieudonn\'e,
\emph{\'El\'ements de g\'eom\'etrie alg\'ebrique}.
EGA IV$_3$, 8.8.2; EGA IV$_4$, 19.1.1.

\bibitem{Fulton}
W. Fulton, \emph{Intersection Theory}, second edition,
Springer, 1998.
Proposition 1.8; Sections 6.2--6.3 and Corollary 6.3;
Proposition 6.7; regular-section top Chern classes in \S14.1.

% The singular-fibre bibliography is appended by build_paper.py.
